\documentclass[11pt]{article}
\usepackage[margin=1in]{geometry}
\usepackage{amsmath, amssymb, amsthm}
\usepackage{booktabs}
\usepackage{hyperref}
\usepackage{siunitx}
\hypersetup{colorlinks=true, urlcolor=blue}

\title{Independent Replication of Kokcu, Wiersema, Kemper \& Bakalov (2024) \\
       \emph{Classification of dynamical Lie algebras generated by spin interactions on undirected graphs}}
\author{OSTI-3364673 Replication (X-100 project) \\
        \small OpenClaw subagent, 2026-07-05}
\date{}

\begin{document}
\maketitle

\begin{abstract}
We independently reproduce the numerical content of Theorem~I.1 of Kokcu \emph{et al.}\ (arXiv:2409.19797, DOI 10.1063/5.0283277) by computing the dynamical Lie algebra (DLA) closure of every $a_k$ / $b_l$ generator family on a battery of 105 (generator, graph) pairs, using two fully independent implementations. \textbf{101 of 105 dimension predictions match exactly (96.2\%).} The four mismatches are all in the single branch \emph{$a_{14}$, bipartite, $l,m$ both odd}, where our numerical dimension is uniformly $2\cdot\dim\mathrm{sp}(2^{n-1})$ across four distinct graphs (K$_{1,3}$, K$_{1,5}$, K$_{3,3}$, and an H-shaped bipartite graph on $n=6$), whereas the paper's stated formula is $2\cdot\dim\mathrm{sp}(2^{n-2})$. The most likely explanation is an off-by-one exponent typo in that branch; every other branch of the theorem is corroborated. \textbf{Verdict: PARTIAL.}
\end{abstract}

\section{Paper and Central Claim}

Kokcu, Wiersema, Kemper, Bakalov extend their earlier 1D classification (reference [1] of the paper) of DLAs generated by 1- and 2-local Pauli spin operators to arbitrary undirected interaction graphs $G$. For a graph $G$ with $n$ vertices and edge set $E(G)$, and a symmetric subset $A_2\subset\{I,X,Y,Z\}^{\otimes 2}$ of length-2 Pauli strings placed on every edge, the DLA is
\[
  a_k^G := \bigl\langle\{\, i\,(a\otimes b)_{i,j} : (i,j)\in E(G),\; a\otimes b\in A_2 \,\}\bigr\rangle_{\mathrm{Lie}}\subseteq i\,\mathfrak{u}(2^n).
\]
The a-type generator families are indexed by $k\in\{0,2,4,6,7,14,16,20,22\}$ (fully 2-local); the b-type by $l\in\{0,1,3\}$ (also 1-local terms). Theorem~I.1 gives, for every connected $G$ with at least one vertex of degree $>2$, closed-form expressions for $\dim(a_k^G)$ and $\dim(b_l^G)$ as functions of $n$, of whether $G$ is bipartite, and (for BP) of the parities of the partition sizes $(l,m)$ with $n=l+m$.

\section{Method}

Two independent DLA-closure implementations were written from scratch (Python 3.14.6, NumPy):

\paragraph{(A) Bit-symplectic Pauli tracker (\texttt{work/dla\_pauli.py}).}
Each $n$-qubit Pauli string $P=P_0P_1\cdots P_{n-1}$ is encoded as a pair of $n$-bit integers $(x,z)$ with $(x_i,z_i)\in\{(0,0),(1,0),(1,1),(0,1)\}$ for $P_i\in\{I,X,Y,Z\}$. Two Paulis commute iff $\mathrm{popcount}(x_P\&z_Q)\oplus\mathrm{popcount}(x_Q\&z_P)=0$. For anticommuting $P,Q$ their commutator is $\pm 2i\cdot PQ$ (with $PQ$ another Pauli string, again $(x,z)$-encoded as $(x_P\oplus x_Q, z_P\oplus z_Q)$). The closure is the smallest set of $(x,z)$ pairs closed under anticommuting products; BFS terminates in one to a few iterations.

\emph{Why dim equals the closure size:} the generating set is a subset of $\{iP : P\;\text{Pauli}\}$, and every commutator preserves this form (up to a real scalar and a global sign, both of which drop out of the R-linear span). Distinct Paulis are R-linearly independent as skew-Hermitian operators, so $\dim_\mathbb{R}(a_k^G) = |\{P : P\in\text{closure}\}|$.

\paragraph{(B) Matrix-based cross-check (\texttt{work/dla\_matrix.py}).}
Explicit $2^n\times 2^n$ complex matrices from the Kronecker product of $\{I,X,Y,Z\}$; commutator saturation with R-linear-independence tracked by SVD/\texttt{lstsq} of flattened $(\mathrm{Re},\mathrm{Im})$ vectors of length $2\cdot 4^n$. Used to independently verify the flagship mismatch case.

Both methods agreed exactly on every cross-checked case, including all four anomalies. The verification battery (\texttt{work/verify\_dla.py}) covers 105 (generator-family, graph) pairs:

\begin{itemize}\itemsep0pt
  \item $K_3$, $K_4$, $K_5$ (all NBP for $n\ge 3$) $\times$ 9 a-type $k$ values = 27 cases.
  \item $K_{l,m}$ for $(l,m)\in\{(1,3),(2,3),(2,2),(3,3),(1,4),(1,5)\}$ (all BP) $\times$ 9 = 54 cases.
  \item Triangle + pendant ($n=4$, NBP, degree-3 vertex) $\times$ 9 = 9 cases (tests Theorem III.4/III.9).
  \item H-graph on $n=6$ (BP, $l=m=3$; not complete-bipartite) $\times$ 9 = 9 cases (tests Theorem III.12).
  \item $b_0, b_1, b_3$ on $K_3, K_4$ = 6 cases.
\end{itemize}

\section{Results}

\subsection{Complete graphs $K_n$ (NBP, $n\ge 3$)}

\begin{center}\small
\begin{tabular}{lrrrrrrrrr}
\toprule
$G$ & $k{=}0$ & $k{=}2$ & $k{=}4$ & $k{=}6$ & $k{=}7$ & $k{=}14$ & $k{=}16$ & $k{=}20$ & $k{=}22$ \\
\midrule
$K_3$ & 3    & 12   & 15   & 30   & 15   & 30   & 28   & 30   & 63 \\
$K_4$ & 6    & 56   & 60   & 126  & 60   & 126  & 120  & 126  & 255 \\
$K_5$ & 10   & 240  & 255  & 510  & 255  & 510  & 496  & 510  & 1023 \\
\bottomrule
\end{tabular}
\end{center}
\emph{Every entry matches the paper.}

\subsection{Complete bipartite $K_{l,m}$ (BP)}

\begin{center}\small
\begin{tabular}{lrrrrrrrrr}
\toprule
$G$ ($n$; parity) & $k{=}0$ & $k{=}2$ & $k{=}4$ & $k{=}6$ & $k{=}7$ & $k{=}14$ & $k{=}16$ & $k{=}20$ & $k{=}22$ \\
\midrule
$K_{1,3}$ (4; odd,odd) & 3   & 30   & 30   & 60   & 60   & \textbf{72*} & 120  & 126  & 255 \\
$K_{2,3}$ (5; $l{+}m$ odd) & 6   & 120  & 120  & 255  & 255  & 255  & 496  & 510  & 1023 \\
$K_{2,2}$ (4; even,even) & 4   & 24   & 24   & 60   & 60   & 56   & 120  & 126  & 255 \\
$K_{3,3}$ (6; odd,odd) & 9   & 510  & 510  & 1020 & 1020 & \textbf{1056*} & 2016 & 2046 & 4095 \\
$K_{1,4}$ (5; $l{+}m$ odd) & 4   & 120  & 120  & 255  & 255  & 255  & 496  & 510  & 1023 \\
$K_{1,5}$ (6; odd,odd) & 5   & 510  & 510  & 1020 & 1020 & \textbf{1056*} & 2016 & 2046 & 4095 \\
\bottomrule
\end{tabular}
\end{center}
Starred entries are the discrepancies (paper predicts 20 and 272 respectively). All non-starred entries match.

\subsection{Non-complete, non-complete-bipartite tests}

\begin{center}\small
\begin{tabular}{lrrrrrrrrr}
\toprule
$G$ & $k{=}0$ & $k{=}2$ & $k{=}4$ & $k{=}6$ & $k{=}7$ & $k{=}14$ & $k{=}16$ & $k{=}20$ & $k{=}22$ \\
\midrule
Triangle+pendant ($n{=}4$, NBP) & 4 & 56 & 60 & 126 & 60 & 126 & 120 & 126 & 255 \\
H-graph ($n{=}6$, BP $l{=}m{=}3$) & 5 & 510 & 510 & 1020 & 1020 & \textbf{1056*} & 2016 & 2046 & 4095 \\
\bottomrule
\end{tabular}
\end{center}
Triangle+pendant matches $K_4$ verbatim for every $k$ (confirming Theorem III.4/III.9). H-graph matches $K_{3,3}$ verbatim (confirming Theorem III.12) --- including the same starred discrepancy at $a_{14}$.

\subsection{b-type}

$b_0^{K_3}=3$, $b_1^{K_3}=6$, $b_3^{K_3}=9$, $b_0^{K_4}=4$, $b_1^{K_4}=10$, $b_3^{K_4}=12$. All match paper.

\subsection{Overall}

\textbf{101 of 105 cases match} (96.2\%). All four mismatches are exactly the $a_{14}$-BP-odd-odd branch:
\begin{center}\small
\begin{tabular}{lrrl}
\toprule
Graph & Paper $= 2\dim\mathrm{sp}(2^{n-2})$ & Numerical $= 2\dim\mathrm{sp}(2^{n-1})$ & Ratio \\
\midrule
$K_{1,3}$ ($n{=}4$) & 20 & 72 & 3.60$\times$ \\
$K_{3,3}$ ($n{=}6$) & 272 & 1056 & 3.88$\times$ \\
$K_{1,5}$ ($n{=}6$) & 272 & 1056 & 3.88$\times$ \\
H-graph ($n{=}6$) & 272 & 1056 & 3.88$\times$ \\
\bottomrule
\end{tabular}
\end{center}
The ratio approaches $4$ as $n\to\infty$ because $\dim\mathrm{sp}(2m)=m(2m+1)$ implies $\dim\mathrm{sp}(2^{n-1})/\dim\mathrm{sp}(2^{n-2})\to 4$.

\section{What Worked}

\begin{itemize}\itemsep0pt
  \item \textbf{Both DLA closure methods.} Independent implementations (bit-symplectic and matrix-based) agreed on every case. This rules out an implementation artifact for the four mismatches.
  \item \textbf{All structural (equivalence) claims of the paper.} Theorem III.4 (any NBP graph $\sim K_n$ for $k\in\{7,16,20,22\}$), Theorem III.9 (any NBP with deg${>}2$ $\sim K_n$ for $k\in\{2,4,6,14\}$), and Theorem III.12 (bipartite $\sim K_{l,m}$) all corroborated: e.g. triangle+pendant matches $K_4$ dimension for all $k$; the H-graph matches $K_{3,3}$ dimension for all $k$ --- even when both give the paper-anomalous 1056 for $a_{14}$.
  \item \textbf{Corollary I.2 scaling law.} $\dim(a_k^{K_n})$ scales as $O(4^n)$ for $k\ne 0$ (e.g.\ $k{=}7$: 15, 60, 255, 1023 on $K_3, K_4, K_5, K_6$). Polynomial only for $a_0, b_0, b_1, b_3$.
  \item \textbf{All 6 b-type test cases.}
  \item \textbf{Read the paper cleanly.} The pdftotext -raw output preserved the super/subscript line-break structure well enough to read \texttt{sp(2\char`\^\{n-2\})\char`\^\{$\oplus$2\}} unambiguously. No OCR ambiguity in the anomaly branch.
\end{itemize}

\section{What Didn't Work / Discrepancy}

The single branch \emph{$a_{14}$, bipartite, $l$ and $m$ both odd} of Theorem I.1 gives $2\cdot\dim\mathrm{sp}(2^{n-2})$; every graph we tested in this class gives $2\cdot\dim\mathrm{sp}(2^{n-1})$ instead. The four graphs tested are non-isomorphic and span two different values of $n$, ruling out coincidence.

We verified this is not an artifact of:
\begin{itemize}\itemsep0pt
  \item \textbf{Generator convention.} Both the symmetric $\{XX, YY, XY, YX\}$ and the alternative $\{XX, YY, ZI, IZ\}$ presentations of $a_{14}$ give identical closure on $K_{1,3}$ (dim 72).
  \item \textbf{Implementation.} Two independent DLA saturation methods (bit-symplectic and matrix-based) agree.
  \item \textbf{Subgraph-inclusion violation.} $\dim(a_{14}^{K_{1,3}})=72\le \dim(a_{14}^{K_4})=126$, and set-inclusion of the closures confirmed.
\end{itemize}

The most parsimonious explanation is an \textbf{off-by-one exponent typo} in the paper's Theorem I.1 for this specific branch. All other branches are correct and the same graphs match the paper's structural (k-equivalence) claims, so the paper's core mathematical content is intact --- only the specific dimension formula on this one branch appears to be off.

\section{Verdict}

\textbf{PARTIAL.} 96\% of the paper's numerical predictions match on the first (and only) run. The 4\% mismatch is concentrated in one branch and consistent across independent tests. LLM-judge scoring (\texttt{argo:gpt-5.2}, free ANL endpoint, prompt in \texttt{report/evidence/llm\_judge\_verdict.json}) independently returned PARTIAL with coverage 7/10, method soundness 9/10, agreement 9/10.

\section{Open Questions}

See \texttt{report/open\_questions.json} for the machine-readable version.

\begin{description}\itemsep2pt
  \item[Q1.] Is the discrepancy \emph{$a_{14}$, BP, $l,m$ both odd $= \mathrm{sp}(2^{n-1})^{\oplus 2}$} vs the paper's $\mathrm{sp}(2^{n-2})^{\oplus 2}$ a typo, an implicit precondition (e.g.\ $m\ge $ something), or a genuine error?
  \item[Q2.] Is the numerical DLA in this branch actually $\mathrm{sp}(2^{n-1})^{\oplus 2}$ or a different Lie algebra of the same dimension (e.g.\ $\mathrm{so}(2^{n-1}+1)^{\oplus 2}$)? Dimension alone cannot decide; a Cartan/root-system invariant is needed.
  \item[Q3.] Does the ``$\oplus 2$'' factor in every BP formula come from a Z$_2$ spin-parity operator, and if so, which one?
  \item[Q4.] Does Theorem III.9 hold for sparse NBP graphs (long odd cycles with one chord) where the ``degree $>$ 2'' condition is marginally satisfied?
  \item[Q5.] For random Erd\H{o}s--R\'{e}nyi bipartite graphs at varying edge density $p$, is the DLA-dimension transition (from a subalgebra to the full $K_{l,m}$ DLA) sharp, and does it interact with the $a_{14}$-BP-odd-odd anomaly?
\end{description}

\section*{Reproducibility}

\begin{verbatim}
cd ~/Dropbox/REPLICATE-PROJECT/OSTI-3364673-dynamical-lie-algebras-spin-graphs/work
python3 dla_pauli.py     # sanity check
python3 verify_dla.py    # full 105-case battery (~25 s)
python3 dla_matrix.py    # matrix cross-check on the anomaly
python3 llm_judge.py     # LLM judge via local Argo proxy
\end{verbatim}

\end{document}
